%% kapa: exam=1 round=34 year=2023 period=2 form=A subject=appraisal_law subject_ko=감정평가관계법규 seq=26 number=26
%% source: Q-Net 감정평가사 기출문제 (cst00309) artlSeq=5211789 file=raw/1008/5211789/2023년도 제34회 감정평가사 1차시험 2교시.pdf
%% conversion: kordoc 4.24.51 -> tools/gen1.py (auto); review=proofread-sonnet; flags=none
\begin{question}{26}
건축법령상 조정(調停) 및 재정(裁定)에 관한 설명으로 옳지 않은 것은?
\begin{choices}
  \item 조정 및 재정을 하기 위하여 국토교통부에 건축분쟁전문위원회를 둔다.
  \item 부득이한 사정으로 연장되지 않는 한 건축분쟁전문위원회는 당사자의 조정신청을 받으면 60일 이내에 절차를 마쳐야 한다.
  \item 조정안을 제시받은 당사자는 제시를 받은 날부터 30일 이내에 수락 여부를 조정위원회에 알려야 한다.
  \item 조정위원회는 필요하다고 인정하면 당사자나 참고인을 조정위원회에 출석하게 하여 의견을 들을 수 있다.
  \item 건축분쟁전문위원회는 재정신청이 된 사건을 조정에 회부하는 것이 적합하다고 인정하면 직권으로 직접 조정할 수 있다.
\end{choices}
\end{question}
\answer{3}
